\documentclass[conference]{IEEEtran}

\usepackage{cite}
\usepackage{amsmath,amssymb,amsfonts}
\usepackage{algorithmic}
\usepackage{graphicx}
\usepackage{textcomp}
\usepackage{xcolor}
\usepackage{booktabs}
\usepackage{multirow}
\usepackage{hyperref}
\usepackage{array}

\begin{document}

\title{Detection and Classification of Peripheral Blood Cells in Microscopic Images using Deep Learning}

\author{
\IEEEauthorblockN{Mohammad Hossain, N. M. Muntasirul Haque, Rafiur Rahman Mashrafi}
\IEEEauthorblockA{Department of Computer Science and Engineering\\
North South University, Dhaka, Bangladesh\\
mohammad.hossain29@northsouth.edu, muntasirul.haque@northsouth.edu, rafiur.mashrafi@northsouth.edu}
}

\maketitle

\begin{abstract}
Peripheral blood smear microscopy is an important diagnostic tool for hematological screening, but manual cell inspection is time-consuming and sensitive to inter-observer variation. In this work, we propose a two-stage deep learning pipeline for peripheral blood cell analysis. First, an object detector localizes cells in TXL-PBC microscopic images. The detected bounding boxes are then used to crop individual cells and build a compact classification dataset. We evaluate two detection models (YOLOv8m and RT-DETR) using standard detection metrics (mAP, precision, and recall), and we compare multiple lightweight and transformer-based classifiers on the cropped three-class dataset (RBC, WBC, and Platelets). Experimental results show that YOLOv8m achieves the best detection performance, while Inception-ResNet-v2 provides the strongest overall classification performance on the test split. The proposed pipeline demonstrates an efficient approach to automate blood cell detection and classification using deep learning.
\end{abstract}

\begin{IEEEkeywords}
Peripheral blood cells, TXL-PBC, object detection, YOLOv8, RT-DETR, cell cropping, image classification, Inception-ResNet-v2, confusion matrix, deep learning.
\end{IEEEkeywords}

\section{Introduction}

Microscopic examination of peripheral blood smears is a routine and clinically valuable procedure because it reveals cell morphology and relative abundance that are essential for hematological screening. However, manual inspection is time-consuming, requires experienced personnel, and is prone to inter-observer variation when slides differ in staining, illumination, focus, or when cells overlap. These practical limitations motivate computer-aided analysis that can localize cells reliably in full-resolution microscope images and then recognize cell types based on morphology.

Recent progress in convolutional neural networks (CNNs) and vision transformers has enabled strong performance in medical image recognition, but robust blood smear analysis still benefits from a workflow that separates localization from recognition and reports metrics that remain informative under class imbalance. In this work, we study a two-stage pipeline on the TXL-PBC dataset: an object detector first localizes RBC, WBC, and platelets, and the resulting bounding boxes are used to build cell-centric crops for downstream classification. This design keeps the system modular and allows a controlled comparison of detectors and classification backbones under the same split and preprocessing setup.

\begin{enumerate}
    \item Choudhary et al. (2025) combine U-Net segmentation with an Inception-ResNet-v2 classifier. By isolating cell regions before classification, their method reduces background influence and highlights the importance of preprocessing and model selection for microscopy images.
    
    \item Sazak and Kotan (2025) evaluate a YOLO-based detector (YOLOv5s) for localizing RBC, WBC, and platelets and report results using precision, recall, and mAP. Their findings emphasize that practical processing choices and data balancing can improve detection stability across varied smear conditions.
    
    \item Wang et al. (2025) explore transformer-based object detection and show that models such as RT-DETR can provide competitive accuracy with end-to-end training. Such approaches are promising for complex scenes where cells overlap, appear at multiple scales, or exhibit non-uniform contrast.
    
    \item A separate line of work combines data augmentation (including GAN-driven synthesis) with region-based detectors such as Mask R-CNN. The key insight is that enriched training variability may improve generalization, although the overall training pipeline becomes more complex.
    
    \item Makki and Al-Dulaimi (2025) discuss hybrid strategies that detect candidate cells using YOLO and then classify cropped instances with a CNN. This two-stage strategy aligns with clinical workflows and motivates using cropping to reduce clutter in the classifier input.
\end{enumerate}

Building on these directions, our project focuses on an efficient detection-to-cropping-to-classification workflow. We evaluate two detectors (YOLOv8m and RT-DETR) for cell localization, and we benchmark five diverse classifiers on a cropped three-class dataset (RBC, WBC, and platelets). To reflect performance on minority classes, we report macro-averaged precision, recall, and F1 in addition to accuracy.

The remainder of the paper is organized as follows. Section II describes the proposed system, dataset, preprocessing, and the models and training setup. Section III reports experimental results and discusses detection and classification performance. Section IV concludes the paper and outlines future work.

\section{Proposed System}

Our proposed system is a two-stage pipeline designed for peripheral blood smear microscopy (Fig.~\ref{fig:pipeline}). In Stage 1, an object detector is trained on full-resolution images to localize RBC, WBC, and platelets by predicting their class labels and bounding boxes. In Stage 2, the bounding boxes are used to extract cell-centric crops, producing a compact classification dataset where the cell occupies most of the frame and surrounding background is minimized. Multiple classification backbones are then fine-tuned on these crops under a consistent protocol to study how different architectures behave when the input is standardized.

This modular design provides practical advantages for both development and evaluation. Detection and classification can be improved independently, and the cropping step reduces scale variation and visual clutter that often complicate direct image-level classification. Moreover, using the same cropped dataset across all classifiers enables a fair comparison across lightweight CNNs and transformer-based models, while macro-averaged metrics highlight minority-class behavior.

\begin{figure}[htbp]
\centerline{\includegraphics[width=\columnwidth]{image1.png}}
\caption{Proposed system pipeline for blood cell detection and classification.}
\label{fig:pipeline}
\end{figure}

\subsection{Dataset Description}

We use the TXL-PBC dataset, which provides microscopic images annotated with bounding boxes for three categories: RBC, WBC, and Platelets. The dataset includes train/validation/test splits and YOLO-format label files. The class distribution is imbalanced, with RBC being the majority class. Table~\ref{tab:distribution} shows the instance-level distribution used in our experiments.

\begin{table}[htbp]
\caption{Instance-Level Class Distribution in TXL-PBC}
\label{tab:distribution}
\centering
\begin{tabular}{l r}
\toprule
\textbf{Class} & \textbf{Instances} \\
\midrule
RBC & 22,412 \\
WBC & 1,806 \\
Platelets & 764 \\
\bottomrule
\end{tabular}
\end{table}

\begin{figure}[htbp]
\centerline{\includegraphics[width=\columnwidth]{image2.png}}
\caption{Dataset statistics and bounding-box distribution (TXL-PBC).}
\label{fig:dataset_stats}
\end{figure}

\subsection{Preprocessing and Dataset Preparation}

To improve visual quality while keeping annotation integrity, we apply a lightweight enhancement pipeline before training. First, a median filter is used to suppress impulsive noise and small artifacts while preserving edges around cell boundaries. Next, Contrast Limited Adaptive Histogram Equalization (CLAHE) is applied to enhance local contrast in a controlled manner, which helps reveal morphology-relevant cues (e.g., boundary sharpness, cytoplasm texture, and nucleus appearance) under varying illumination and staining. Importantly, we do not perform geometric resizing or warping on the full microscope images during this step so that the original YOLO-format bounding boxes remain valid and aligned with the images.

For the classification stage, we generate cell-level crops using the available bounding-box annotations. Each crop is saved under its corresponding class label and later resized to the model input resolution ($224 \times 224$) during training. This strategy focuses the classifier on the cell of interest, reduces background variation, and makes comparisons across classification backbones more reliable.

\begin{figure}[htbp]
\centerline{\includegraphics[width=\columnwidth]{image3.png}}
\caption{Sample images from the raw TXL-PBC dataset.}
\label{fig:raw_samples}
\end{figure}

\begin{figure}[htbp]
\centerline{\includegraphics[width=\columnwidth]{image4.png}}
\caption{Raw samples with ground-truth bounding boxes.}
\label{fig:raw_bbox}
\end{figure}

\begin{figure}[htbp]
\centerline{\includegraphics[width=\columnwidth]{image5.png}}
\caption{Preprocessing evidence (before vs. after Median + CLAHE).}
\label{fig:preprocessing}
\end{figure}

\begin{figure}[htbp]
\centerline{\includegraphics[width=\columnwidth]{image6.png}}
\caption{Samples from the preprocessed dataset.}
\label{fig:preprocessed_samples}
\end{figure}

Using the annotation boxes, we crop each cell region and store the resulting patches in three class folders (Platelets, RBC, and WBC). Compared with classifying full microscope images, cropping isolates the relevant morphology and reduces the influence of unrelated regions (e.g., empty background or nearby cells). The same cropped dataset is then used across all classification backbones to ensure consistent training conditions and a fair evaluation.

\begin{figure}[htbp]
\centerline{\includegraphics[width=\columnwidth]{image7.png}}
\caption{Cropped classification dataset samples (3 classes $\times$ 6 examples).}
\label{fig:cropped_samples}
\end{figure}

\subsection{Models and Training Setup}

\subsubsection{Object Detection Models}

We evaluate two modern detectors: YOLOv8m and RT-DETR. YOLOv8m is a one-stage detector optimized for fast inference and strong accuracy, while RT-DETR is a transformer-based detector designed for real-time performance with end-to-end training. Both models are trained on the same dataset split with comparable augmentation policies to ensure a meaningful comparison.

For YOLOv8m, we train for 30 epochs with image size 640 and batch size 16 (seed=42, workers=2). For RT-DETR, we train for 30 epochs with image size 640 and batch size 8 (reduced to fit GPU memory). We apply standard geometric and photometric augmentations including small rotations, translations, scaling, and color jitter. During inference, non-maximum suppression (NMS) is used to remove duplicate detections.

\subsubsection{Cell Classification Models}

For classification, we compare five architectures covering both lightweight CNNs and transformer-based networks: Inception-ResNet-v2, GhostNet, FocalNet, MaxViT, and PVTv2. Each model is trained on the cropped three-class dataset using the same data split and preprocessing pipeline.

\textbf{Inception-ResNet-v2:} Inception-ResNet-v2 combines the multi-scale feature extraction of Inception blocks with residual connections. This design stabilizes deep training and captures both fine-grained texture and broader shape patterns, which are important for distinguishing WBC morphology from RBC and platelets. In our experiments, it achieved the best overall macro-F1 on the test set.

\textbf{GhostNet:} GhostNet targets efficiency by generating more feature maps from inexpensive operations, reducing redundancy in convolutions. It is suitable for resource-constrained environments while maintaining competitive accuracy. On the cropped cell dataset, GhostNet performed very close to the best model, indicating that lightweight CNNs can still be strong baselines for medical classification when images are well-cropped.

\textbf{FocalNet:} FocalNet introduces focal modulation to aggregate contextual information efficiently, allowing the model to capture both local and global patterns without the full cost of attention. This is beneficial for blood cell images where texture and boundary cues matter. The model achieved high accuracy and robust per-class performance in our evaluation.

\textbf{MaxViT:} MaxViT combines window-based self-attention with grid attention to balance local detail and global context. It is effective for images where objects occupy different spatial regions and where contextual cues help classification. In our tests, MaxViT achieved strong overall performance with stable per-class precision and recall.

\textbf{PVTv2:} PVTv2 is a pyramid vision transformer that builds multi-scale representations similarly to CNN feature pyramids. Although it is powerful for general vision tasks, its performance on our cropped dataset was lower than the other models, suggesting that careful tuning and larger training data may be required for transformer-only backbones in this setting.

All classifiers are trained with input size $224 \times 224$ for 30 epochs using batch size 32. We use the AdamW optimizer (lr=$3 \times 10^{-4}$, weight decay=$1 \times 10^{-4}$) and fix the random seed (seed=42) to improve reproducibility. Training-time augmentation includes random horizontal flips and small rotations, followed by normalization.

\section{Results and Discussion}

\subsection{Detection Results}

Table~\ref{tab:detection} reports detection performance on both the validation and test splits. Both YOLOv8m and RT-DETR achieve high precision and recall, indicating that the models localize most cells while producing relatively few false positives. On the test split, YOLOv8m achieves the best overall performance (mAP50-95 = 0.895), while RT-DETR remains competitive (mAP50-95 = 0.866). Overall, these results suggest that a strong one-stage detector provides an excellent accuracy-speed trade-off for smear images, and transformer-based detection can be a viable alternative when computational budget allows.

\begin{table}[htbp]
\caption{Detection Performance Comparison (Val and Test)}
\label{tab:detection}
\centering
\begin{tabular}{l l c c c c}
\toprule
\textbf{Model} & \textbf{Split} & \textbf{P} & \textbf{R} & \textbf{mAP50} & \textbf{mAP50--95} \\
\midrule
\textbf{YOLOv8m} & \textbf{Val} & \textbf{0.953} & \textbf{0.962} & \textbf{0.979} & \textbf{0.881} \\
RT-DETR & Val & 0.927 & 0.961 & 0.970 & 0.858 \\
\textbf{YOLOv8m} & \textbf{Test} & \textbf{0.950} & \textbf{0.962} & \textbf{0.977} & \textbf{0.895} \\
RT-DETR & Test & 0.941 & 0.963 & 0.976 & 0.866 \\
\bottomrule
\end{tabular}
\end{table}

Beyond aggregate scores, Fig.~\ref{fig:confusion_det} (normalized confusion matrix) summarizes class-wise detection behavior for YOLOv8m. Most predictions lie on the diagonal, which indicates that detected objects are generally assigned the correct class. The remaining errors are more common for visually small or low-contrast instances, where platelets may be missed or confused with tiny RBC fragments, a challenge amplified by the instance-level class imbalance in TXL-PBC.

As shown in Fig.~\ref{fig:f1_curve}, the F1-score remains strong across a broad range of confidence thresholds, which is useful for selecting a stable operating point at inference time. Fig.~\ref{fig:pr_curve} presents class-wise precision-recall curves that stay near the upper-right region, reflecting strong precision even at high recall. Finally, qualitative examples in Fig.~\ref{fig:qualitative_det} confirm that YOLOv8m produces tight bounding boxes across different cell densities and remains robust when multiple adjacent cells appear in the same field.

\begin{figure}[htbp]
\centerline{\includegraphics[width=\columnwidth]{image8.png}}
\caption{Normalized confusion matrix for YOLOv8m detection.}
\label{fig:confusion_det}
\end{figure}

\begin{figure}[htbp]
\centerline{\includegraphics[width=\columnwidth]{image9.png}}
\caption{F1-score vs. confidence threshold (YOLOv8m).}
\label{fig:f1_curve}
\end{figure}

\begin{figure}[htbp]
\centerline{\includegraphics[width=\columnwidth]{image10.png}}
\caption{Precision--Recall curve by class (YOLOv8m).}
\label{fig:pr_curve}
\end{figure}

\begin{figure}[htbp]
\centerline{\includegraphics[width=\columnwidth]{image11.png}}
\caption{Qualitative detection results (YOLOv8m) on sample images.}
\label{fig:qualitative_det}
\end{figure}

\subsection{Classification Results}

After generating the cropped classification dataset, we evaluate five backbones: Inception-ResNet-v2, GhostNet, FocalNet, MaxViT, and PVT-v2. All models are fine-tuned using the same data split, input resolution, optimizer, and augmentation policy to ensure a controlled comparison. Because the dataset is imbalanced at the instance level, we report macro-averaged precision, recall, and F1 in addition to accuracy so that minority-class performance is reflected more fairly.

Table~\ref{tab:classification} summarizes classification performance on the test split. Inception-ResNet-v2 achieves the highest Macro-F1 (0.9953) with near-perfect accuracy, indicating that the cropped cell-centric representation is highly separable for the three classes. GhostNet and FocalNet perform very close to the best model, suggesting that carefully designed lightweight architectures can be highly effective when the input is well-cropped. MaxViT remains strong but shows a small drop in macro-recall, which suggests slightly higher sensitivity to minority-class instances. In contrast, PVT-v2 exhibits notably lower macro-recall and Macro-F1, implying that transformer-only backbones may require additional tuning, regularization, or more training data to match the top-performing models in this setting.

\begin{table}[htbp]
\caption{Overall Classification Performance on the Test Set}
\label{tab:classification}
\centering
\begin{tabular}{l c c c c}
\toprule
\textbf{Model} & \textbf{Acc.} & \textbf{Macro-P} & \textbf{Macro-R} & \textbf{Macro-F1} \\
\midrule
\textbf{Inception-ResNet-v2} & \textbf{0.9995} & \textbf{0.9975} & \textbf{0.9932} & \textbf{0.9953} \\
GhostNet & 0.9989 & 0.9931 & 0.9973 & 0.9952 \\
FocalNet & 0.9989 & 0.9951 & 0.9930 & 0.9940 \\
MaxViT & 0.9984 & 0.9966 & 0.9898 & 0.9932 \\
PVT-v2 & 0.9830 & 0.9156 & 0.8018 & 0.8555 \\
\bottomrule
\end{tabular}
\end{table}

As shown in Fig.~\ref{fig:macro_f1}, the Macro-F1 comparison indicates that the top three models are closely matched. To understand class-wise behavior for the best classifier, Fig.~\ref{fig:confusion_cls} (normalized confusion matrix) and Fig.~\ref{fig:per_class} (per-class precision/recall/F1) show that most errors are concentrated in minority classes, which is expected given the smaller number of platelet and WBC samples. Despite this imbalance, the per-class metrics remain high, indicating strong generalization beyond the majority RBC class.

\begin{figure}[htbp]
\centerline{\includegraphics[width=\columnwidth]{image13.png}}
\caption{Macro-F1 comparison of classification models on the test set.}
\label{fig:macro_f1}
\end{figure}

\begin{figure}[htbp]
\centerline{\includegraphics[width=\columnwidth]{image14.png}}
\caption{Normalized confusion matrix for the best classifier (Inception-ResNet-v2).}
\label{fig:confusion_cls}
\end{figure}

\begin{figure}[htbp]
\centerline{\includegraphics[width=\columnwidth]{image15.png}}
\caption{Per-class precision, recall, and F1-score (Inception-ResNet-v2).}
\label{fig:per_class}
\end{figure}

The one-vs-rest precision-recall and ROC curves (Fig.~\ref{fig:pr_ovr} and Fig.~\ref{fig:roc_ovr}) further confirm strong separability for each class. In addition, sample test-set predictions in Fig.~\ref{fig:sample_pred} provide qualitative evidence that the model remains reliable across diverse appearances in the test split. Overall, these findings support the effectiveness of bounding-box based cropping: by centering the cell and reducing surrounding clutter, it simplifies the classification task and enables reliable benchmarking across different network families.

\begin{figure}[htbp]
\centerline{\includegraphics[width=\columnwidth]{image16.png}}
\caption{One-vs-rest Precision--Recall curves (Inception-ResNet-v2).}
\label{fig:pr_ovr}
\end{figure}

\begin{figure}[htbp]
\centerline{\includegraphics[width=\columnwidth]{image17.png}}
\caption{One-vs-rest ROC curves (Inception-ResNet-v2).}
\label{fig:roc_ovr}
\end{figure}

\begin{figure}[htbp]
\centerline{\includegraphics[width=\columnwidth]{image18.png}}
\caption{Sample test predictions of the best classifier (Inception-ResNet-v2).}
\label{fig:sample_pred}
\end{figure}

Overall, bounding-box based cropping simplifies classification by standardizing the object scale and reducing background variability, which is particularly important in smear images where multiple cells may appear in the same field. The remaining mistakes tend to occur in minority classes, underscoring that class imbalance can still influence errors even when crops are clean. For broader deployment, it would be valuable to evaluate cross-domain generalization across different laboratories and staining protocols and to investigate techniques such as stain normalization or domain adaptation to reduce distribution shift.

\section{Conclusion and Future Work}

In this work, we developed and evaluated a two-stage deep learning pipeline for peripheral blood cell analysis on TXL-PBC microscopy images. The pipeline combines (i) object detection for localizing RBC, WBC, and platelets in full-resolution images with (ii) bounding-box based cropping to construct cell-centric inputs for classification, followed by (iii) benchmarking of multiple classification backbones under a consistent training protocol. Experimentally, YOLOv8m achieved the best detection performance on the test split (mAP50-95 = 0.895), while Inception-ResNet-v2 delivered the strongest classification results (Accuracy = 0.9995, Macro-F1 = 0.9953). These results demonstrate that a modular detection-to-classification workflow can provide accurate and efficient automation for peripheral blood smear analysis.

Future work will focus on improving robustness and clinical applicability. We plan to extend the pipeline to additional blood cell subtypes and to evaluate generalization across different laboratories, scanners, and staining protocols. Incorporating stain normalization or domain adaptation may reduce distribution shift, and analyzing difficult cases (e.g., overlapping cells or low-contrast platelets) can guide targeted improvements. We also aim to study end-to-end or multi-task training that jointly optimizes detection and classification, and to add calibration or uncertainty estimation so the system can better communicate confidence in ambiguous cases.

\begin{thebibliography}{10}

\bibitem{txlpbc}
TXL-PBC Dataset, GitHub repository (lugan113/TXL-PBC\_Dataset).

\bibitem{inception}
C. Szegedy et al., ``Inception-v4, Inception-ResNet and the Impact of Residual Connections on Learning,'' 2016.

\bibitem{yolov8}
Ultralytics, ``YOLOv8: Real-Time Object Detection,'' 2023.

\bibitem{rtdetr}
Y. Lv et al., ``RT-DETR: Real-Time Detection Transformer,'' 2023.

\bibitem{ghostnet}
H. Han et al., ``GhostNet: More Features from Cheap Operations,'' 2020.

\bibitem{focalnet}
J. Yang et al., ``FocalNet: Focal Modulation Networks,'' 2022.

\bibitem{maxvit}
Z. Tu et al., ``MaxViT: Multi-Axis Vision Transformer,'' 2022.

\bibitem{pvtv2}
W. Wang et al., ``PVTv2: Improved Baselines with Pyramid Vision Transformer,'' 2022.

\bibitem{makki}
S. Makki and A. Al-Dulaimi, ``Peripheral Blood Cell Classification from Microscopic Images: A Review,'' Diagnostics, 15(1):22, 2025. doi:10.3390/diagnostics15010022.

\bibitem{additional}
Additional recent works summarized in the literature review (doi:10.1038/s41598-025-91720-7; doi:10.1186/s12911-025-03027-2).

\end{thebibliography}

\end{document}
